%=====================================================================
% Datos · Casos normales
%=====================================================================
\subsection*{Casos normales}
\addcontentsline{toc}{subsection}{Casos normales}

Un caso normal es el que un caso de uso describe en su flujo normal, sin alternativas ni excepciones. La mayoría de los registros lo son: cuentas registradas y activas que juegan partidas clasificatorias y compran objetos. Dos ejemplos bastan para ver todo lo que un caso normal deja escrito.

\subsubsection*{Una partida clasificatoria completa}
\addcontentsline{toc}{subsubsection}{Una partida clasificatoria completa}

La partida 1 enfrenta a \at{luis\_quo} y \at{marta\_muros} en el modo clásico, con reloj de cinco minutos. Esta consulta reúne la partida, sus dos participaciones y las monedas que cada uno recibió al terminar:

\consulta{normalpartida}
% Archivo generado por bd/exportar_datos.py. No editar a mano.
\begin{center}\scriptsize\setlength{\tabcolsep}{2pt}
\begin{longtable}{|L{46.3pt}|L{57.2pt}|L{42.1pt}|L{62.4pt}|L{59.2pt}|L{41.0pt}|L{29.8pt}|L{38.2pt}|L{55.2pt}|}
\hline
\textbf{\hspace{0pt}Partida} & \textbf{\hspace{0pt}Modo} & \textbf{\hspace{0pt}Forma} & \textbf{\hspace{0pt}Jugador} & \textbf{\hspace{0pt}Resultado} & \textbf{\hspace{0pt}Elo inicial} & \textbf{\hspace{0pt}Elo final} & \textbf{\hspace{0pt}Diferencia} & \textbf{\hspace{0pt}Monedas} \\ \hline
\endfirsthead
\hline
\textbf{\hspace{0pt}Partida} & \textbf{\hspace{0pt}Modo} & \textbf{\hspace{0pt}Forma} & \textbf{\hspace{0pt}Jugador} & \textbf{\hspace{0pt}Resultado} & \textbf{\hspace{0pt}Elo inicial} & \textbf{\hspace{0pt}Elo final} & \textbf{\hspace{0pt}Diferencia} & \textbf{\hspace{0pt}Monedas} \\ \hline
\endhead
1 & CLASICO & META & luis\_\allowbreak{}quo & GANADA & 1000 & 1016 & 16 & 100 \\ \hline
1 & CLASICO & META & marta\_\allowbreak{}muros & PERDIDA & 1000 & 984 & -16 & 25 \\ \hline
\end{longtable}
\end{center}


Los dos empezaron con el elo inicial de 1000 (\mbox{CU-02} \mbox{RN-06}). Con elos iguales, la victoria vale 16 puntos y la derrota los mismos 16 en contra, y la diferencia la calcula la columna \at{diferencia\_elo}. Cada participante recibe un movimiento de monedas de tipo \at{PARTIDA}: 100 el ganador y 25 el perdedor (\mbox{CU-25} \mbox{RN-08}). Las quince jugadas de la partida alternan los turnos, mezclan movimientos y muros y terminan con luis en la fila 9, la meta del jugador que sale de la fila 1:

\consulta{normaljugadas}
% Archivo generado por bd/exportar_datos.py. No editar a mano.
\begin{center}\scriptsize\setlength{\tabcolsep}{3pt}
\begin{longtable}{|L{0.052\linewidth}|L{0.120\linewidth}|L{0.109\linewidth}|L{0.052\linewidth}|L{0.061\linewidth}|L{0.044\linewidth}|L{0.109\linewidth}|L{0.055\linewidth}|L{0.251\linewidth}|}
\hline
\textbf{Número} & \textbf{Jugador} & \textbf{Tipo} & \textbf{Origen} & \textbf{Destino} & \textbf{Surco} & \textbf{Orientación} & \textbf{Tiempo (ms)} & \textbf{Fecha} \\ \hline
\endfirsthead
\hline
\textbf{Número} & \textbf{Jugador} & \textbf{Tipo} & \textbf{Origen} & \textbf{Destino} & \textbf{Surco} & \textbf{Orientación} & \textbf{Tiempo (ms)} & \textbf{Fecha} \\ \hline
\endhead
1 & luis\_\allowbreak{}quo & MOVIMIENTO & e1 & e2 & \textit{NULL} & \textit{NULL} & 7148 & 2026-07-05 17:00:07.148 \\ \hline
2 & marta\_\allowbreak{}muros & MURO & \textit{NULL} & \textit{NULL} & a7 & HORIZONTAL & 6067 & 2026-07-05 17:00:13.215 \\ \hline
3 & luis\_\allowbreak{}quo & MOVIMIENTO & e2 & e3 & \textit{NULL} & \textit{NULL} & 4986 & 2026-07-05 17:00:18.201 \\ \hline
4 & marta\_\allowbreak{}muros & MOVIMIENTO & e9 & d9 & \textit{NULL} & \textit{NULL} & 3905 & 2026-07-05 17:00:22.106 \\ \hline
5 & luis\_\allowbreak{}quo & MOVIMIENTO & e3 & e4 & \textit{NULL} & \textit{NULL} & 2824 & 2026-07-05 17:00:24.930 \\ \hline
6 & marta\_\allowbreak{}muros & MURO & \textit{NULL} & \textit{NULL} & g7 & HORIZONTAL & 10743 & 2026-07-05 17:00:35.673 \\ \hline
7 & luis\_\allowbreak{}quo & MOVIMIENTO & e4 & e5 & \textit{NULL} & \textit{NULL} & 9662 & 2026-07-05 17:00:45.335 \\ \hline
8 & marta\_\allowbreak{}muros & MOVIMIENTO & d9 & d8 & \textit{NULL} & \textit{NULL} & 8581 & 2026-07-05 17:00:53.916 \\ \hline
9 & luis\_\allowbreak{}quo & MOVIMIENTO & e5 & e6 & \textit{NULL} & \textit{NULL} & 7500 & 2026-07-05 17:01:01.416 \\ \hline
10 & marta\_\allowbreak{}muros & MURO & \textit{NULL} & \textit{NULL} & b3 & VERTICAL & 6419 & 2026-07-05 17:01:07.835 \\ \hline
11 & luis\_\allowbreak{}quo & MOVIMIENTO & e6 & e7 & \textit{NULL} & \textit{NULL} & 5338 & 2026-07-05 17:01:13.173 \\ \hline
12 & marta\_\allowbreak{}muros & MOVIMIENTO & d8 & d7 & \textit{NULL} & \textit{NULL} & 4257 & 2026-07-05 17:01:17.430 \\ \hline
13 & luis\_\allowbreak{}quo & MOVIMIENTO & e7 & e8 & \textit{NULL} & \textit{NULL} & 3176 & 2026-07-05 17:01:20.606 \\ \hline
14 & marta\_\allowbreak{}muros & MURO & \textit{NULL} & \textit{NULL} & h2 & HORIZONTAL & 11095 & 2026-07-05 17:01:31.701 \\ \hline
15 & luis\_\allowbreak{}quo & MOVIMIENTO & e8 & e9 & \textit{NULL} & \textit{NULL} & 10014 & 2026-07-05 17:01:41.715 \\ \hline
\end{longtable}
\end{center}


En un movimiento, la jugada tiene casilla de origen y de destino y no tiene surco; en un muro, al revés. Es la regla de subtipo que hace cumplir la restricción \at{CK\_Jugada\_tipo}. La fecha de cada jugada lleva milisegundos, y el tiempo consumido es el que se descuenta del reloj del jugador.

\subsubsection*{Compras en la tienda}
\addcontentsline{toc}{subsubsection}{Compras en la tienda}

Una compra deja dos registros en la misma transacción: un movimiento de monedas de tipo \at{COMPRA}, con importe negativo, y el objeto en el inventario de la cuenta, con origen \at{COMPRA} (\mbox{CU-38} \mbox{RN-08}). La consulta los une y comprueba además si el objeto quedó equipado:

\consulta{normalcompra}
% Archivo generado por bd/exportar_datos.py. No editar a mano.
\begin{center}\scriptsize\setlength{\tabcolsep}{2pt}
\begin{longtable}{|L{44.4pt}|L{54.9pt}|L{92.8pt}|L{43.9pt}|L{43.3pt}|L{74.4pt}|L{50.8pt}|L{31.3pt}|}
\hline
\textbf{\hspace{0pt}Movimiento} & \textbf{\hspace{0pt}Usuario} & \textbf{\hspace{0pt}Objeto} & \textbf{\hspace{0pt}Rareza} & \textbf{\hspace{0pt}Importe} & \textbf{\hspace{0pt}Fecha} & \textbf{\hspace{0pt}Origen en UsuarioObjeto} & \textbf{\hspace{0pt}Equipado} \\ \hline
\endfirsthead
\hline
\textbf{\hspace{0pt}Movimiento} & \textbf{\hspace{0pt}Usuario} & \textbf{\hspace{0pt}Objeto} & \textbf{\hspace{0pt}Rareza} & \textbf{\hspace{0pt}Importe} & \textbf{\hspace{0pt}Fecha} & \textbf{\hspace{0pt}Origen en UsuarioObjeto} & \textbf{\hspace{0pt}Equipado} \\ \hline
\endhead
10 & luis\_\allowbreak{}quo & EMOTE\_\allowbreak{}RISA & RARO & -120 & 2026-07-21 18:00 & COMPRA & sí \\ \hline
13 & luis\_\allowbreak{}quo & PEON\_\allowbreak{}CRISTAL & RARO & -250 & 2026-07-27 18:00 & COMPRA & sí \\ \hline
21 & marta\_\allowbreak{}muros & MURO\_\allowbreak{}LADRILLO & COMUN & -150 & 2026-08-16 19:00 & COMPRA & sí \\ \hline
23 & pedro\_\allowbreak{}peon & EMOTE\_\allowbreak{}RISA & RARO & -120 & 2026-08-21 17:00 & COMPRA & sí \\ \hline
\end{longtable}
\end{center}


El importe de cada compra es el precio del objeto en el catálogo, y no el que envía el cliente (\mbox{CU-38} \mbox{RN-03}). Los cuatro objetos comprados quedaron equipados en su ranura (\mbox{CU-14}).
